\chapter{2006年全国II卷（理）}
\section{单选题}
\begin{problem}
已知集合 \(M = \{x \mid x < 3\}\)，\(N = \{x \mid \log_2 x > 1\}\)，则 \(M \cap N =\)（\quad）

\choices
    {\(\varnothing\)}
    {\(\{x \mid 0 < x < 3\}\)}
    {\(\{x \mid 1 < x < 3\}\)}
    {\(\{x \mid 2 < x < 3\}\)}
\end{problem}

\begin{answer}
D
\end{answer}

\begin{solution}
由 \(\log_2x>1\) 且底数 \(2>1\)，得 \(x>2\)，所以 \(N=\{x\mid x>2\}\). 因此
\(M\cap N=\{x\mid 2<x<3\}\)，故选 D.
\end{solution}

\begin{problem}
函数 \(y = \sin 2x \cos 2x\) 的最小正周期是（\quad）

\choices
    {\(2\pi\)}
    {\(4\pi\)}
    {\(\frac{\pi}{4}\)}
    {\(\frac{\pi}{2}\)}
\end{problem}

\begin{answer}
D
\end{answer}

\begin{solution}
利用倍角公式，\(y=\sin2x\cos2x=\frac12\sin4x\). 因为 \(\sin4x\) 的最小正周期为 \(\frac{2\pi}{4}=\frac\pi2\)，所以该函数的最小正周期为 \(\frac\pi2\)，故选 D.
\end{solution}

\begin{problem}
\(\frac{3}{(1 - \i)^2} =\)（\quad）

\choices
    {\(\frac{3}{2}\i\)}
    {\(-\frac{3}{2}\i\)}
    {\(1\)}
    {\(-1\)}
\end{problem}

\begin{answer}
A
\end{answer}

\begin{solution}
\((1-\i)^2=1-2\i+\i^2=-2\i\)，从而
\(\frac{3}{(1-\i)^2}=\frac{3}{-2\i}=\frac{3\i}{2}\)，故选 A.
\end{solution}

\begin{problem}
过球的一条半径的中点，作垂直于该半径的平面，则所得截面的面积与球的表面积的比为（\quad）

\choices
    {\(\frac{3}{16}\)}
    {\(\frac{9}{16}\)}
    {\(\frac{3}{8}\)}
    {\(\frac{9}{32}\)}
\end{problem}

\begin{answer}
A
\end{answer}

\begin{solution}
设球的半径为 \(R\). 截面所在平面到球心的距离为半径中点到球心的距离 \(\frac R2\)，所以截面圆半径的平方为
\(R^2-\left(\frac R2\right)^2=\frac{3R^2}{4}\). 截面面积与球表面积之比为
\(\frac{\pi\cdot\frac{3R^2}{4}}{4\pi R^2}=\frac{3}{16}\)，故选 A.
\end{solution}

\begin{problem}
已知 \(\triangle ABC\) 的顶点 \(B\)，\(C\) 在椭圆 \(\frac{x^2}{3} + y^2 = 1\) 上，顶点 \(A\) 是椭圆的一个焦点，且椭圆的另外一个焦点在 \(BC\) 边上，则 \(\triangle ABC\) 的周长是（\quad）

\choices
    {\(2\sqrt{3}\)}
    {\(6\)}
    {\(4\sqrt{3}\)}
    {\(12\)}
\end{problem}

\begin{answer}
C
\end{answer}

\begin{solution}
设椭圆的两个焦点为 \(A\)，\(F\)，则长半轴为 \(\sqrt3\)，故椭圆上一点到两焦点的距离之和为 \(2\sqrt3\). 因为 \(F\) 在边 \(BC\) 上，所以 \(BC=BF+FC\). 于是三角形周长为
\(AB+BC+AC=(AB+BF)+(AC+FC)=2\sqrt3+2\sqrt3=4\sqrt3\)，故选 C.
\end{solution}

\begin{problem}
已知函数 \(f(x) = \ln x + 1\)（\(x > 0\)），则 \(f(x)\) 的反函数为（\quad）

\choices
    {\(y = \e^{x + 1}\)（\(x \in \R\)）}
    {\(y = \e^{x - 1}\)（\(x \in \R\)）}
    {\(y = \e^{x + 1}\)（\(x > 1\)）}
    {\(y = \e^{x - 1}\)（\(x > 1\)）}
\end{problem}

\begin{answer}
B
\end{answer}

\begin{solution}
由 \(y=\ln x+1\) 解得 \(x=\e^{y-1}\). 原函数的值域为 \(\R\)，交换自变量与因变量后，反函数为 \(y=\e^{x-1}\)，其定义域为 \(\R\)，故选 B.
\end{solution}

\begin{problem}
如图，平面 \(\alpha \perp\) 平面 \(\beta\)，\(A \in \alpha\)，\(B \in \beta\)，\(AB\) 与两平面 \(\alpha\)，\(\beta\) 所成的角分别为 \(\frac{\pi}{4}\) 和 \(\frac{\pi}{6}\). 过 \(A\)，\(B\) 分别作两平面交线的垂线，垂足为 \(A'\)，\(B'\)，则 \(AB:A'B'=\)（\quad）

\bitmapfigure[width=0.42\linewidth]{img/2006/national_paper_2_science/q7_fig1.png}

\choices
    {\(2:1\)}
    {\(3:1\)}
    {\(3:2\)}
    {\(4:3\)}
\end{problem}

\begin{answer}
A
\end{answer}

\begin{solution}
设两平面的交线为 \(x\) 轴，取 \(\alpha\) 为 \(z=0\)，\(\beta\) 为 \(y=0\). 设 \(AB=L\)，并写成
\(A=(u_1,v,0)\)，\(B=(u_2,0,w)\). 直线 \(AB\) 与平面 \(\alpha\)，\(\beta\) 所成角分别为 \(\frac\pi4\)，\(\frac\pi6\)，所以
\(\lvert w\rvert=L\sin\frac\pi4\)，\(\lvert v\rvert=L\sin\frac\pi6\). 由于 \(A'=(u_1,0,0)\)，\(B'=(u_2,0,0)\)，有
\[
 A'B'^2=(u_2-u_1)^2=L^2-v^2-w^2=L^2-\frac{L^2}{4}-\frac{L^2}{2}=\frac{L^2}{4}
\]
因此 \(A'B'=\frac L2\)，即 \(AB:A'B'=2:1\)，故选 A.
\end{solution}

\begin{problem}
函数 \(y = f(x)\) 的图象与函数 \(g(x) = \log_2 x\)（\(x > 0\)）的图象关于原点对称，则 \(f(x)\) 的表达式为（\quad）

\choices
    {\(f(x) = \frac{1}{\log_2 x}\)（\(x > 0\)）}
    {\(f(x) = \frac{1}{\log_2(-x)}\)（\(x < 0\)）}
    {\(f(x) = -\log_2 x\)（\(x > 0\)）}
    {\(f(x) = -\log_2(-x)\)（\(x < 0\)）}
\end{problem}

\begin{answer}
D
\end{answer}

\begin{solution}
设 \((x,y)\) 在 \(y=f(x)\) 的图象上. 关于原点的对称点 \((-x,-y)\) 在 \(y=\log_2x\) 的图象上，因此 \(-x>0\) 且
\(-y=\log_2(-x)\). 所以 \(x<0\) 时 \(f(x)=-\log_2(-x)\)，故选 D.
\end{solution}

\begin{problem}
已知双曲线 \(\frac{x^2}{a^2} - \frac{y^2}{b^2} = 1\) 的一条渐近线方程为 \(y = \frac{4}{3} x\)，则双曲线的离心率为（\quad）

\choices
    {\(\frac{5}{3}\)}
    {\(\frac{4}{3}\)}
    {\(\frac{5}{4}\)}
    {\(\frac{3}{2}\)}
\end{problem}

\begin{answer}
A
\end{answer}

\begin{solution}
双曲线的渐近线为 \(y=\pm\frac ba x\)，故 \(\frac ba=\frac43\). 又双曲线满足 \(c^2=a^2+b^2\)，所以
\[
 e=\frac ca=\sqrt{1+\frac{b^2}{a^2}}=\sqrt{1+\frac{16}{9}}=\frac53
\]
故选 A.
\end{solution}

\begin{problem}
若 \(f(\sin x) = 3 - \cos 2x\)，则 \(f(\cos x) =\)（\quad）

\choices
    {\(3 - \cos 2x\)}
    {\(3 - \sin 2x\)}
    {\(3 + \cos 2x\)}
    {\(3 + \sin 2x\)}
\end{problem}

\begin{answer}
C
\end{answer}

\begin{solution}
由 \(\cos2x=1-2\sin^2x\)，得
\(f(\sin x)=3-\cos2x=2+2\sin^2x\). 由于 \(\sin x\) 可以取遍 \([-1,1]\)，故对 \(t\in[-1,1]\) 有 \(f(t)=2+2t^2\). 因而
\(f(\cos x)=2+2\cos^2x=3+\cos2x\)，故选 C.
\end{solution}

\begin{problem}
设 \(S_{n}\) 是等差数列 \(\{a_{n}\}\) 的前 \(n\) 项和，若 \(\frac{S_{3}}{S_{6}} = \frac{1}{3}\)，则 \(\frac{S_{6}}{S_{12}} =\)（\quad）

\choices
    {\(\frac{3}{10}\)}
    {\(\frac{1}{3}\)}
    {\(\frac{1}{8}\)}
    {\(\frac{1}{9}\)}
\end{problem}

\begin{answer}
A
\end{answer}

\begin{solution}
设等差数列的首项为 \(a_1\)，公差为 \(d\). 由求和公式，
\(S_3=3(a_1+d)\)，\(S_6=3(2a_1+5d)\). 由 \(\frac{S_3}{S_6}=\frac13\)，得
\(3(a_1+d)=2a_1+5d\)，即 \(a_1=2d\). 由于题设比值有意义，\(S_6\ne0\)，所以由
\(S_6=3(4d+5d)=27d\) 得 \(d\ne0\). 又 \(S_{12}=6(4d+11d)=90d\)，从而
\(\frac{S_6}{S_{12}}=\frac{27}{90}=\frac3{10}\)，故选 A.
\end{solution}

\begin{problem}
函数 \(\sum_{n=1}^{19} |x - n|\) 的最小值为（\quad）

\choices
    {\(190\)}
    {\(171\)}
    {\(90\)}
    {\(45\)}
\end{problem}

\begin{answer}
C
\end{answer}

\begin{solution}
记 \(F(x)=\sum_{n=1}^{19}\lvert x-n\rvert\). 对 \(j=1\)，\(2\)，\(\dots\)，\(9\)，由三角不等式有
\(\lvert x-j\rvert+\lvert x-(20-j)\rvert\geq 20-2j\)，并且 \(\lvert x-10\rvert\geq0\). 因此
\[
 F(x)\geq(18+16+\cdots+2)+0=90
\]
当且仅当 \(x=10\) 时，上述各项同时取等号，所以最小值为 \(90\)，故选 C.
\end{solution}

\section{填空题}
\begin{problem}
在 \(\left(x^{4} + \frac{1}{x}\right)^{10}\) 的展开式中常数项是\fillinblank{}.
\end{problem}

\begin{answer}
45
\end{answer}

\begin{solution}
通项中从 \(x^4\) 取 \(k\) 次、从 \(\frac1x\) 取 \(10-k\) 次，所得项为
\(\mathrm C_{10}^{k}(x^4)^k\left(\frac1x\right)^{10-k}=\mathrm C_{10}^{k}x^{5k-10}\). 要成为常数项，需满足 \(5k-10=0\)，即 \(k=2\). 因此常数项为 \(\mathrm C_{10}^{2}=45\).
\end{solution}

\begin{problem}
已知 \(\triangle ABC\) 的三个内角 \(A\)，\(B\)，\(C\) 成等差数列，且 \(AB = 1\)，\(BC = 4\)，则边 \(BC\) 上的中线 \(AD\) 的长为\fillinblank{}.
\end{problem}

\begin{answer}
\(\sqrt{3}\)
\end{answer}

\begin{solution}
因为 \(A\)，\(B\)，\(C\) 成等差数列，所以 \(A+C=2B\). 又 \(A+B+C=\pi\)，故 \(B=\frac\pi3\). 在三角形 \(ABC\) 中，\(AB=1\)，\(BC=4\)，由余弦定理得
\(AC^2=AB^2+BC^2-2\cdot AB\cdot BC\cos B=1+16-4=13\).
设 \(D\) 为 \(BC\) 的中点，则 \(BD=2\). 由中线的平方公式，
\(AB^2+AC^2=2(AD^2+BD^2)\)，所以 \(AD^2=\frac{1+13}{2}-4=3\)，从而 \(AD=\sqrt3\).
\end{solution}

\begin{problem}
过点 \((1, \sqrt{2})\) 的直线 \(l\) 将圆 \((x - 2)^2 + y^2 = 4\) 分成两段弧，当劣弧所对的圆心角最小时，直线 \(l\) 的斜率 \(k =\fillinblank{}\).
\end{problem}

\begin{answer}
\(\frac{\sqrt{2}}{2}\)
\end{answer}

\begin{solution}
圆心为 \(O=(2,0)\)，半径为 \(2\)，而给定点 \(P=(1,\sqrt2)\) 满足 \(OP=\sqrt3<2\)，在圆内. 对于过圆内一点的弦，劣弧所对的圆心角越小，弦长越短；过 \(P\) 的最短弦垂直于 \(OP\)，所以直线 \(l\perp OP\). 直线 \(OP\) 的斜率为 \(-\sqrt2\)，故 \(k=\frac{1}{\sqrt2}=\frac{\sqrt2}{2}\).
\end{solution}

\begin{problem}
一个社会调查机构就某地居民的月收入调查了\(10000\text{ 人}\)，并根据所得数据画了样本的频率分布直方图（如图）. 为了分析居民的收入与年龄，学历，职业等方面的关系，要从这\(10000\text{ 人}\)中再用分层抽样方法抽出\(100\text{ 人}\)作进一步调查，则在月收入\([2500, 3000)\text{ 元}\)段应抽出\fillinblank{}\(\text{ 人}\).

\bitmapfigure[width=0.78\linewidth]{img/2006/national_paper_2_science/q16_fig1.png}
\end{problem}

\begin{answer}
25
\end{answer}

\begin{solution}
由直方图可知，月收入 \([2500,3000)\) 元这一组的频率密度为 \(0.0005\)，组距为 \(500\text{ 元}\)，所以该组频率为 \(0.0005\times500=0.25\). 该组有 \(10000\times0.25=2500\text{ 人}\). 按比例从这一层抽取的人数为 \(100\times\frac{2500}{10000}=25\text{ 人}\).
\end{solution}

\section{解答题}
\begin{problem}
已知向量 \(\bs{a} = (\sin \theta, \sqrt{3})\)，\(\bs{b} = (1, \cos \theta)\)，\(-\frac{\pi}{2} < \theta < \frac{\pi}{2}\).

\begin{enumerate}
\item 若 \(\bs{a} \perp \bs{b}\)，求 \(\theta\)；
\item 求 \(|\bs{a} + \bs{b}|\) 的最大值.
\end{enumerate}
\end{problem}

\begin{solution}
\begin{enumerate}
\item 由 \(\bs{a}\perp\bs{b}\)，得 \(\bs{a}\cdot\bs{b}=\sin\theta+\sqrt3\cos\theta=0\). 因为 \(-\frac\pi2<\theta<\frac\pi2\)，所以 \(\cos\theta>0\)，从而 \(\tan\theta=-\sqrt3\). 因此 \(\theta=-\frac\pi3\).
\item
\[
\lvert\bs{a}+\bs{b}\rvert^2=(1+\sin\theta)^2+(\sqrt3+\cos\theta)^2=5+2(\sin\theta+\sqrt3\cos\theta)=5+4\cos\left(\theta-\frac\pi6\right)\leq 9
\]
当 \(\theta=\frac\pi6\) 时等号成立，且该值在给定范围内，所以 \(|\bs{a}+\bs{b}|\) 的最大值为 \(3\).
\end{enumerate}
\end{solution}

\begin{problem}
某批产品成箱包装，每箱\(5\text{ 件}\)，一用户在购进该批产品前先取出\(3\text{ 箱}\)，再从每箱中任意取出\(2\text{ 件}\)产品进行检验. 设取出的第一，二，三箱中分别有\(0\text{ 件}\)，\(1\text{ 件}\)，\(2\text{ 件}\)二等品，其余为一等品.

\begin{enumerate}
\item 用 \(\xi\) 表示抽检的\(6\text{ 件}\)产品中二等品的件数，求 \(\xi\) 的分布列及 \(\xi\) 的数学期望；
\item 若抽检的\(6\text{ 件}\)产品中有\(2\text{ 件}\)或\(2\text{ 件以上}\)二等品，用户就拒绝购买这批产品，求这批产品被用户拒绝购买的概率.
\end{enumerate}
\end{problem}

\begin{solution}
\begin{enumerate}
\item 第一箱没有二等品，设从第二、三箱抽出的二等品数分别为 \(\xi_2\)，\(\xi_3\)，则 \(\xi=\xi_2+\xi_3\). 从第二箱抽取 \(2\) 件，得到 \(P(\xi_2=0)=\frac{\mathrm C_4^2}{\mathrm C_5^2}=\frac35\)，\(P(\xi_2=1)=\frac{\mathrm C_1^1\mathrm C_4^1}{\mathrm C_5^2}=\frac25\). 从第三箱抽取 \(2\) 件，得到 \(P(\xi_3=0)=\frac{\mathrm C_3^2}{\mathrm C_5^2}=\frac3{10}\)，\(P(\xi_3=1)=\frac{\mathrm C_2^1\mathrm C_3^1}{\mathrm C_5^2}=\frac35\)，\(P(\xi_3=2)=\frac{\mathrm C_2^2}{\mathrm C_5^2}=\frac1{10}\). 两箱抽取相互独立，因此
\(P(\xi=0)=\frac35\cdot\frac3{10}=\frac9{50}\)，
\(P(\xi=1)=\frac35\cdot\frac35+\frac25\cdot\frac3{10}=\frac{24}{50}\)，
\(P(\xi=2)=\frac35\cdot\frac1{10}+\frac25\cdot\frac35=\frac{15}{50}\)，
\(P(\xi=3)=\frac25\cdot\frac1{10}=\frac2{50}\). 所以分布列为
\begin{center}
\begin{tabular}{c|cccc}
\(\xi\) & \(0\) & \(1\) & \(2\) & \(3\) \\
\hline
\(P(\xi)\) & \(\frac9{50}\) & \(\frac{24}{50}\) & \(\frac{15}{50}\) & \(\frac2{50}\)
\end{tabular}
\end{center}
因此 \(E(\xi)=0\cdot\frac9{50}+1\cdot\frac{24}{50}+2\cdot\frac{15}{50}+3\cdot\frac2{50}=\frac65\).
\item 用户拒绝购买的事件为 \(\{\xi\geq2\}\)，所以拒绝购买的概率为 \(P(\xi\geq2)=P(\xi=2)+P(\xi=3)=\frac{15}{50}+\frac2{50}=\frac{17}{50}\).
\end{enumerate}
\end{solution}

\begin{problem}
如图，在直三棱柱 \(ABC - A_{1}B_{1}C_{1}\) 中，\(AB = BC\)，\(D\)，\(E\) 分别为 \(BB_{1}\)，\(AC_{1}\) 的中点.

\begin{enumerate}
\item 证明：\(ED\) 为异面直线 \(BB_{1}\) 与 \(AC_{1}\) 的公垂线；
\item 设 \(AA_{1} = AC = \sqrt{2}AB\)，求二面角 \(A_{1} - AD - C_{1}\) 的大小.
\end{enumerate}

\bitmapfigure[width=0.38\linewidth]{img/2006/national_paper_2_science/q19_fig1.png}
\end{problem}

\begin{solution}
\begin{enumerate}
\item 设 \(O\) 为 \(AC\) 的中点. 在三角形 \(ACC_1\) 中，\(O\)，\(E\) 分别为 \(AC\)，\(AC_1\) 的中点，所以 \(OE\parallel CC_1\)，且 \(OE=\frac12CC_1\). 又 \(D\) 为 \(BB_1\) 的中点，且直棱柱的侧棱平行且相等，所以 \(BD\parallel CC_1\)，且 \(BD=\frac12CC_1\). 因而四边形 \(BOED\) 为平行四边形，得到 \(ED\parallel BO\).

由于 \(AB=BC\)，且 \(O\) 是 \(AC\) 的中点，所以 \(BO\perp AC\). 设 \(O_1\) 为 \(A_1C_1\) 的中点，则 \(OO_1\parallel CC_1\). 又侧棱 \(CC_1\perp\) 底面 \(ABC\)，故 \(BO\perp CC_1\)，从而 \(BO\perp OO_1\). 直线 \(AC\) 与 \(OO_1\) 在 \(O\) 相交且都在平面 \(ACC_1A_1\) 内，因此 \(BO\perp\) 平面 \(ACC_1A_1\)，从而 \(ED\perp AC_1\). 同时 \(ED\parallel BO\)，而 \(BB_1\perp\) 底面 \(ABC\)，故 \(ED\perp BB_1\). 直线 \(ED\) 分别与 \(AC_1\)，\(BB_1\) 在 \(E\)，\(D\) 相交，所以 \(ED\) 是这两条异面直线的公垂线.
\item 设 \(AB=s>0\). 由题设 \(AA_1=AC=\sqrt2s\)，而 \(AB=BC=s\)，故 \(AC^2=AB^2+BC^2\)，底面三角形在 \(B\) 处为直角. 取 \(s=1\) 建立空间直角坐标系：\(A=(1,0,0)\)，\(B=(0,0,0)\)，\(C=(0,1,0)\)，\(A_1=(1,0,\sqrt2)\)，\(C_1=(0,1,\sqrt2)\)，\(D=(0,0,\frac{\sqrt2}{2})\).

平面 \(A_1AD\) 的方程为 \(y=0\)，可取法向量 \(\bs{n}_1=(0,1,0)\). 平面 \(ADC_1\) 内有方向向量 \(\overrightarrow{AD}=(-1,0,\frac{\sqrt2}{2})\) 和 \(\overrightarrow{AC_1}=(-1,1,\sqrt2)\)，可取法向量 \(\bs{n}_2=(1,-1,\sqrt2)\)，因为 \(\bs{n}_2\cdot\overrightarrow{AD}=0\)，\(\bs{n}_2\cdot\overrightarrow{AC_1}=0\). 因而二面角的平面角 \(\varphi\) 满足
\(\cos\varphi=\frac{|\bs{n}_1\cdot\bs{n}_2|}{|\bs{n}_1||\bs{n}_2|}=\frac12\). 所以 \(\varphi=\frac\pi3\).
\end{enumerate}
\end{solution}

\begin{problem}
设函数 \(f(x) = (x + 1)\ln (x + 1)\). 若对所有的 \(x \geqslant 0\)，都有 \(f(x) \geqslant ax\) 成立，求实数 \(a\) 的取值范围.
\end{problem}

\begin{solution}
令 \(g(x)=f(x)-ax=(x+1)\ln(x+1)-ax\)，其中 \(x\geq0\). 则 \(g(0)=0\)，且 \(g'(x)=\ln(x+1)+1-a\).

当 \(a\leq1\) 时，\(x\geq0\) 有 \(\ln(x+1)\geq0\)，所以 \(g'(x)\geq1-a\geq0\). 因此 \(g\) 在 \([0,+\infty)\) 上单调不减，从而 \(g(x)\geq g(0)=0\)，即 \(f(x)\geq ax\) 对所有 \(x\geq0\) 成立.

当 \(a>1\) 时，取任意 \(x\) 满足 \(0<x<\e^{a-1}-1\)，则 \(\ln(x+1)<a-1\)，从而在 \([0,x]\) 上 \(g'(t)<0\)，故 \(g(x)<g(0)=0\). 这说明此时不可能对所有 \(x\geq0\) 成立. 因此 \(a\) 的取值范围为 \((-\infty,1]\).
\end{solution}

\begin{problem}
已知抛物线 \(x^{2} = 4y\) 的焦点为 \(F\)，\(A\)，\(B\) 是抛物线上的两动点，且 \(\overrightarrow{AF} = \lambda \overrightarrow{FB}\)（\(\lambda > 0\)）. 过 \(A\)，\(B\) 两点分别作抛物线的切线，设其交点为 \(M\).

\begin{enumerate}
\item 证明 \(\overrightarrow{FM} \cdot \overrightarrow{AB}\) 为定值；
\item 设 \(\triangle ABM\) 的面积为 \(S\)，写出 \(S = f(\lambda)\) 的表达式，并求 \(S\) 的最小值.
\end{enumerate}
\end{problem}

\begin{solution}
\begin{enumerate}
\item 设抛物线上的点 \(A\)，\(B\) 分别对应参数 \(u\)，\(v\)，则
\(A=(2u,u^2)\)，\(B=(2v,v^2)\)，且焦点 \(F=(0,1)\). 由 \(\overrightarrow{AF}=\lambda\overrightarrow{FB}\) 得
\[
(-2u,1-u^2)=\lambda(2v,v^2-1)
\]
由第一坐标，\(u=-\lambda v\)；代入第二坐标，得
\(1-\lambda^2v^2=\lambda(v^2-1)\)，即 \((1+\lambda)(1-\lambda v^2)=0\). 因为 \(\lambda>0\)，所以 \(v^2=\frac1\lambda\)，进而 \(u^2=\lambda\)，且 \(uv=-1\).

参数为 \(u\) 的切线方程为 \(y=ux-u^2\)，参数为 \(v\) 的切线方程为 \(y=vx-v^2\)，联立得切线交点 \(M=(u+v,uv)=(u+v,-1)\). 因此
\(\overrightarrow{FM}=(u+v,-2)\)，\(\overrightarrow{AB}=(2(v-u),(v-u)(u+v))\)，从而
\(\overrightarrow{FM}\cdot\overrightarrow{AB}=2(v-u)(u+v)-2(v-u)(u+v)=0\). 这说明该数量积为定值 \(0\).
\item 由 \(\overrightarrow{AF}=\lambda\overrightarrow{FB}\) 且 \(\lambda>0\)，点 \(F\) 在线段 \(AB\) 上，所以 \(AB=AF+FB\). 又
\(AF=\sqrt{4u^2+(u^2-1)^2}=u^2+1=\lambda+1\)，\(FB=\sqrt{4v^2+(v^2-1)^2}=v^2+1=\frac1\lambda+1\). 因而
\[
AB=\lambda+\frac1\lambda+2=\left(\sqrt\lambda+\frac1{\sqrt\lambda}\right)^2
\]
另外，
\[
FM=\sqrt{(u+v)^2+4}=\sqrt{u^2+v^2+2uv+4}=\sqrt{\lambda+\frac1\lambda+2}=\sqrt\lambda+\frac1{\sqrt\lambda}
\]
由第（1）问，\(FM\perp AB\)，且垂足为 \(F\)，所以
\[
S=\frac12AB\cdot FM=\frac12\left(\sqrt\lambda+\frac1{\sqrt\lambda}\right)^3
\]
由均值不等式，\(\sqrt\lambda+\frac1{\sqrt\lambda}\geq2\)，等号当且仅当 \(\lambda=1\) 时成立. 故 \(S\geq4\)，最小值为 \(4\)，在 \(\lambda=1\) 时取得.
\end{enumerate}
\end{solution}

\begin{problem}
设数列 \(\{a_{n}\}\) 的前 \(n\) 项和为 \(S_{n}\)，且方程 \(x^{2} - a_{n}x - a_{n} = 0\) 有一根为 \(S_{n} - 1\)，\(n = 1\)，\(2\)，\(3\)，\(\dots\).

\begin{enumerate}
\item 求 \(a_{1}\)，\(a_{2}\)；
\item 求 \(\{a_{n}\}\) 的通项公式.
\end{enumerate}
\end{problem}

\begin{solution}
\begin{enumerate}
\item 当 \(n=1\) 时，\(S_1=a_1\)，所以方程的一个根为 \(a_1-1\). 代入方程得
\((a_1-1)^2-a_1(a_1-1)-a_1=0\)，化简为 \(1-2a_1=0\)，故 \(a_1=\frac12\).

当 \(n=2\) 时，\(S_2-1=a_2-\frac12\). 代入方程得
\(\left(a_2-\frac12\right)^2-a_2\left(a_2-\frac12\right)-a_2=0\)，化简为 \(\frac14-\frac32a_2=0\)，故 \(a_2=\frac16\).
\item 由根为 \(S_n-1\)，对任意 \(n\) 有
\[
(S_n-1)^2-a_n(S_n-1)-a_n=0
\]
即 \(S_n^2-2S_n+1-a_nS_n=0\). 当 \(n\geq2\) 时，\(a_n=S_n-S_{n-1}\)，代入得
\(S_n(S_{n-1}-2)+1=0\)，即
\[
S_n=\frac1{2-S_{n-1}}
\]
又 \(S_1=\frac12\). 下面用数学归纳法证明 \(S_n=\frac{n}{n+1}\).

当 \(n=1\) 时，\(S_1=\frac12=\frac1{1+1}\)，结论成立. 假设 \(S_k=\frac{k}{k+1}\)，则由递推式
\(S_{k+1}=\frac1{2-\frac{k}{k+1}}=\frac{k+1}{k+2}\)，所以结论对 \(k+1\) 也成立. 因此对所有正整数 \(n\)，都有 \(S_n=\frac{n}{n+1}\).

于是 \(a_1=S_1=\frac12\)，而当 \(n\geq2\) 时
\[
a_n=S_n-S_{n-1}=\frac{n}{n+1}-\frac{n-1}{n}=\frac1{n(n+1)}
\]
该公式对 \(n=1\) 也成立，故 \(a_n=\frac1{n(n+1)}\)（\(n=1\)，\(2\)，\(3\)，\(\dots\)）.
\end{enumerate}
\end{solution}
